\chapter{Examples, Trading, and Repository Integration}

A monograph about repository integration should not stay at the level of definitions. It
should show how those definitions illuminate the code that users actually see.

\section*{Examples as semantic stress tests}
The examples shipped in the repository exercise several different faces of the system:
proof-oriented examples, tensor and shape examples, mixed-mode examples, zero-change
checking examples, and the trading examples that drive the prompt-based generation story.
These are not merely demos. They are semantic stress tests for the thesis that one local
to-global architecture can support multiple workflows.

\section*{Why the trading system matters}
The trading story is useful because it forces the repository to coordinate many ideas at
once. There is natural-language intent about discovering or constructing a strategy. There
are structural constraints such as position sizing, leverage bounds, exposure rules, and
interface boundaries between market data, strategy kernels, risk engines, execution, and
portfolio state. There are semantic judgments about whether a strategy meaningfully realizes
a stated thesis. There are proof ambitions around risk-critical invariants. And there is
actual code generation and assembly.

The hybrid cover-oriented generation pipeline is therefore not theoretical ornament in the
trading example. It is exactly the mechanism needed to keep multiple modules and multiple
kinds of evidence from drifting apart. A trading project becomes a cover of local sections:
market-data sections, signal sections, risk sections, execution sections, portfolio
sections, compliance sections, and final assembly. The question ``did generation succeed?''
becomes a gluing question.

\begin{figure}[h]
\centering
\begin{tikzpicture}[
  node distance=0.95cm,
  site/.style={draw=chapterblue, rounded corners=2pt, fill=chapterblue!6, minimum width=2.35cm, minimum height=0.8cm, align=center},
  >=Latex
]
\node[site] (market) {market\_data};
\node[site, right=of market] (signal) {signal\_geometry};
\node[site, right=of signal] (kernel) {strategy\_kernel};
\node[site, below=of signal] (risk) {risk\_engine};
\node[site, right=of risk] (exec) {execution};
\node[site, right=of exec] (portfolio) {portfolio};
\node[site, below=of exec] (compliance) {compliance};
\node[site, below=of kernel] (main) {main / global section};
\draw[->, thick, chapterblue] (market) -- (signal);
\draw[->, thick, chapterblue] (signal) -- (kernel);
\draw[->, thick, chapterblue] (kernel) -- (risk);
\draw[->, thick, chapterblue] (risk) -- (exec);
\draw[->, thick, chapterblue] (exec) -- (portfolio);
\draw[->, thick, chapterblue] (portfolio) -- (main);
\draw[->, thick, chapterblue] (compliance) -- (main);
\draw[->, dashed, accentgold] (risk) -- (compliance);
\draw[->, dashed, accentgold] (kernel) -- (main);
\end{tikzpicture}
\caption{A repository-native trading cover. Solid arrows follow the execution and state flow; dashed arrows indicate proof and compliance constraints that must still glue into the final global section.}
\end{figure}

The newer scalar-valued typing layer fits naturally into this trading picture. Different
mathematical lenses---for example stochastic calculus, financial mathematics, or control
theory---can assign different scalar weights to the same type elaboration. Those weights do
not prove the system correct, but they help the ideation stage decide which elaborations
deserve expansion into the expensive multi-module cover.

\section*{Repository-wide integration map}
The repository can be read through the following semantic map.
\begin{itemize}[leftmargin=2em, itemsep=0.5em]
  \item \repo{src/deppy/types/}: the foundational object language.
  \item \repo{src/deppy/hybrid/core/}: layered type semantics, trust, checker, and geometry.
  \item \repo{src/deppy/hybrid/mixed_mode/}: user-facing bridge from Python syntax to obligations.
  \item \repo{src/deppy/hybrid/discharge/} and \repo{lean_translation/}: proof-facing boundary.
  \item \repo{src/deppy/hybrid/type_expansion/}: planning and architectural elaboration.
  \item \repo{src/deppy/hybrid/anti_hallucination/}: structural and semantic artifact validation.
  \item \repo{src/deppy/equivalence/}: local-to-global comparison of programs and predicates.
  \item \repo{src/deppy/render/}: turning typed and trust-bearing results into readable diagnostics.
  \item \repo{src/deppy/theory_packs/} and theory libraries: modular local solver geometries.
  \item \repo{agent/}: the orchestration loop that tries to grow a global section from prompt-local material.
\end{itemize}

Once this map is visible, the repository stops looking like a pile of unrelated prototypes.
It looks like a large, uneven, but conceptually coherent attempt to turn local semantics
into a universal organizing principle.
